\chapter{Architecture styles and trade-offs}

\section{Architecture is a set of consequential decisions}
Architecture describes the significant structures, boundaries, and constraints
that shape a system. A diagram is useful only when it helps someone decide,
implement, operate, or change the system. A box labelled “service” is not an
architecture unless its responsibility, communication, data ownership, and
failure behaviour are understood.

\section{Common styles}
\noindent\begin{tabularx}{\textwidth}{@{}p{0.23\textwidth}p{0.26\textwidth}X@{}}
\toprule
Style & Helps when & Costs and failure modes \\
\midrule
Modular monolith & One team needs strong local transactions and simple deployment & A shared process can become tangled; scaling and isolation are coarse. \\
Layered system & Responsibilities have a stable vertical flow & Layers become pass-through plumbing or leak abstractions. \\
Ports and adapters & Domain logic should remain independent of infrastructure & Too many interfaces obscure a small system; mappings must be maintained. \\
Event-driven & Producers and consumers need temporal decoupling & Ordering, retries, schemas, and debugging become distributed concerns. \\
Microservices & Teams need independent ownership and deployment at meaningful boundaries & Network failure, data ownership, versioning, and operational cost multiply. \\
\bottomrule
\end{tabularx}

None is the universally mature choice. A microservice is not automatically more
scalable, and a monolith is not automatically a prototype. The decision should
be driven by boundaries of change, ownership, failure isolation, deployment
independence, data consistency, and capacity rather than by fashion.

\section{Architecture decision records}
An architecture decision record (ADR) keeps a decision close to its reasons:

\begin{enumerate}
  \item Context: what forces and constraints exist?
  \item Decision: what will we do?
  \item Alternatives: what serious options were considered?
  \item Consequences: what becomes easier, harder, or newly risky?
  \item Status: proposed, accepted, superseded, or rejected.
\end{enumerate}

An ADR is not a vote and not a substitute for measurement. It is a memory aid.
If a decision was made under an assumption about traffic, team size, or a vendor,
record the assumption and a trigger for revisiting it.

\section{Data ownership}
In a distributed architecture, the hardest boundary is often data. If two
services write the same tables, ownership is nominal. If they have separate
stores, cross-service transactions become workflows with compensation or
reconciliation. The right choice depends on consistency requirements and
failure tolerance.

For Ledgerline, keep order state in one database while the domain is small. A
separate payment service may make sense when legal responsibility, risk controls,
or independent scaling require it. Splitting it merely to create more boxes
would make failures harder to recover from without solving a real constraint.

\section{Architecture fitness functions}
A \term{fitness function} is a repeatable check for an architectural property.
Examples:

\begin{itemize}
  \item a dependency rule that prevents the domain package importing HTTP code;
  \item a migration test that proves an old and new application version can coexist;
  \item a load test that measures p95 latency at a stated arrival rate;
  \item a security test that rejects access across a tenant boundary;
  \item an accessibility test that checks keyboard navigation for a critical flow.
\end{itemize}

Fitness functions do not prove the architecture is good. They catch regression
against properties the team has decided matter. A metric without a decision
attached can become decorative or actively misleading.

\section{Architecture review}
Review the smallest set of questions that could change the design:

\begin{itemize}
  \item What are the critical user journeys and failure consequences?
  \item Which state is authoritative and which copies are disposable?
  \item What happens when every dependency is slow, unavailable, or inconsistent?
  \item How is a change released and rolled back while old clients remain?
  \item Which trust boundaries, secrets, personal data, and audit obligations exist?
\end{itemize}

\begin{tradeoffs}
Centralization simplifies coordination and can create a bottleneck. Distribution
can improve isolation and team autonomy and can also multiply failure modes. The
correct architecture is the smallest one that meets the real constraints with
known consequences.
\end{tradeoffs}

\begin{exercise}
Choose a modular monolith, a service split, or an event-driven design for an
application that has one team, 20 requests per second, a strict audit trail, and
one database. Defend the choice, then identify the first signal that would make
you reconsider it.
\end{exercise}

\section{Chapter summary}
Architecture is about boundaries and consequences. Select a style for ownership,
failure, change, data, and capacity reasons. Record assumptions in ADRs and turn
important architectural properties into automated or observable fitness checks.
